% Problemas del capítulo 1. Generado por tools/problemas/gen.py; no editar a mano.
\section{Problemas}\label{sec:problemas-01}

Graduados: \({\star}\) directo, \({\star}{\star}\) para pensar, \({\star}{\star}{\star}\) desafío. Las respuestas están en el apéndice~\ref{ap:respuestas}; conviene intentarlos antes de mirarlas.

\begin{problema}[\({\star}\)]\label{pb:01-1}
Rotá el punto \(p=(4,1)\) un ángulo de \(30^\circ\) en sentido antihorario alrededor del origen. Comprobá que la distancia al origen no cambió.
\end{problema}

\begin{problema}[\({\star}\)]\label{pb:01-2}
Con números complejos: multiplicá \(z=3+4i\) por \(i\). ¿Qué ángulo giró \(z\)? ¿Cambió su módulo?
\end{problema}

\begin{problema}[\({\star}{\star}\)]\label{pb:01-3}
Desde una estación \(S=(1000;\ 2000)\)~m, en coordenadas \((E,N)\), se mide un punto a \(120\)~m con acimut \(\mathrm{Az}=135^\circ\). Calculá sus coordenadas. ¿Cuánto vale el ángulo matemático \(\theta\), medido desde el este en sentido antihorario?
\end{problema}

\begin{problema}[\({\star}{\star}\)]\label{pb:01-4}
Girá \(p=(5,1)\) un ángulo de \(90^\circ\) antihorario, pero alrededor del centro \(c=(2,1)\), no del origen. Hacelo en tres pasos: llevar \(c\) al origen, girar y volver.
\end{problema}

\begin{problema}[\({\star}{\star}{\star}\)]\label{pb:01-5}
Hay instrumentos que miden ángulos con precisión de un segundo de arco. Si el ángulo hacia un punto que está a \(2\)~km cambia \(1''\), ¿cuánto se corre el punto de costado? Y en el espacio: aplicá a \(e_1\) primero un giro de \(90^\circ\) alrededor del eje \(z\) y después uno de \(90^\circ\) alrededor del eje \(x\); repetí en el otro orden. ¿Da lo mismo?
\end{problema}
